\documentclass{article}
\usepackage[T1]{fontenc}
\usepackage{lmodern}
\usepackage{graphicx}
\usepackage{amsmath,amssymb}
\usepackage[a4paper,margin=2cm]{geometry}
\usepackage[absolute,overlay]{textpos}
\setlength{\TPHorizModule}{1mm}
\setlength{\TPVertModule}{1mm}
\usepackage[indonesian]{babel}
\usepackage{hyperref}
\setlength{\emergencystretch}{2em}
% unit: Halaman 5

\begin{document}

% === Halaman 5 (python/native) ===

Contoh 1:  Misalkan n = 6 dan w1 = 1, w2 = 5, w3 = 6,  w4 = 11, w5 = 14, dan w6 = 20.

Plainteks:  111001010110000000011000

Plainteks dibagi menjadi blok yang panjangnya 6, kemudian setiap bit di dalam blok dikalikan 

dengan wi yang berkoresponden sesuai dengan persamaan (1):

 Blok plainteks ke-1 

: 111001 

 Kriptogram 


: (1  1) + (1  5) + (1  6) +  (0  11) + (0 x 14) + (1 x 20)   = 32

 Blok plainteks ke-2 

: 010110 

 Kriptogram 


: (1  5) + (1  11) + (1  14) = 30

 Blok plainteks ke-3 

: 000000 

 Kriptogram 


: 0

 Blok plainteks ke-4 

: 011000 

 Kriptogram 


: (1  5) + (1  6) = 11

Jadi, cipherteks yang dihasilkan:  32  30  0  11

5

\end{document}
